\subsection*{Milton Hernández}
El estudiante \textbf{Milton Hernández} contribuyó al proyecto con las siguientes tareas:

\begin{taskdone}{Adaptación del repositorio del Proyecto 1}
    Creación de fork del repositorio del Proyecto 1 y adaptación al nuevo contexto de trabajo con los nuevos integrantes del grupo.

    \tasknote{Se mantiene el sistema de deportistas desarrollado previamente como base para integrar los nuevos algoritmos.}
\end{taskdone}

\begin{taskdone}{Estructura inicial del informe LaTeX}
    Preparación del documento LaTeX con la estructura base del informe según los requisitos de la tarea.
\end{taskdone}

\begin{taskdone}{Implementación de Quick Sort con 4 variantes}
    Implementación de Quick Sort utilizando el esquema de partición de Lomuto con cuatro estrategias de selección de pivote.
    \begin{itemize}
        \item Creación de rama \texttt{quick\_sort}
        \item Implementación del esquema de partición de Lomuto
        \item Variantes de pivote (Primer elemento, Último elemento, Mediana de tres, Elemento aleatorio)
        \item Análisis de comportamiento en mejor, peor y caso promedio y elección de pivote para implementación final
        \item Merge a rama principal
    \end{itemize}
\end{taskdone}

\begin{taskdone}{Implementación de Quick Select}
    Implementación del algoritmo Quick Select para encontrar el k-ésimo elemento sin ordenar todo el arreglo.
    \begin{itemize}
        \item Creación de rama \texttt{quick\_select}
        \item Reutilización del esquema de partición de Quick Sort
        \item Integración con funcionalidad ``k-ésimo mejor deportista''
        \item Merge a rama principal
    \end{itemize}
\end{taskdone}

\begin{taskdone}{Análisis de complejidad temporal de algoritmos iterativos de ordenamiento}
    Análisis experimental del comportamiento de \texttt{quicksort} y \texttt{mergesort} en el mejor, peor y caso promedio.
    
    Comparación con los de los algoritmos implementados en la tarea anterior.
\end{taskdone}

\begin{taskdone}{Ranking de deportistas}
    Generar un ranking con los mejores $N$ deportistas, utilizando el puntaje como criterio.
\end{taskdone}